\chapter{Обзор предметной области}
Использование оптических технологий там, представляет собой большой прогресс по сравнению с использовавшимися ранее электричекими. Меньшая длина волны позволяет достичь большей частоты сигналов, в частности при передаче данных. Такая частота дает возможность передавать потоки информации в несколько терабит в секунду. Важными преимуществами ВОЛС являются такие факторы, как малое затухание сигналов, позволяющее, при использовании современных технологий, строить участки оптических систем в сто и более километров без ретрансляции, высокая помехозащищенность, связанная с малой восприимчивостью оптического волокна к электромагнитным помехам, и многие другие.

Волоконная оптика - это раздел оптики, изучающий процессы в оптоволокне. Оптическое волокно — нить из оптически прозрачного материала (стекло, пластик), используемая для переноса света внутри себя посредством полного внутреннего отражения.

Для объяснения этого эффекта рассмотрим процессы происходящие при прохожении светом границы двух сред:

\begin{figure}
\includegraphics[width=0.5\textwidth]{img/refraction.pdf}
\caption{Преломление света}
\end{figure}